% preamble/packages.tex — Paketimporte
% ═════════════════════════════════════════════════════════════════════

% ── Sprache & Zeichenkodierung ────────────────────────────────────────
\usepackage[ngerman]{babel}
\usepackage[autostyle=true, german=guillemets]{csquotes}   % dt. Anführungszeichen

% ── Schriften (LuaLaTeX + fontspec) ──────────────────────────────────
% Libertinus: freie OpenType-Familie, in TeX Live 2020+ enthalten
\usepackage{fontspec}
\setmainfont{Libertinus Serif}
\setsansfont{Libertinus Sans}
\setmonofont[Scale=MatchLowercase]{Source Code Pro}         % für Code

\usepackage{unicode-math}
\setmathfont{Libertinus Math}                               % passend zur Serifenschrift

% Mikrotypographie: Randausgleich, Unterschneidung, Wortzwischenräume
\usepackage[
  protrusion = true,
  expansion  = true,
  final,
]{microtype}

% ── Seitenlayout & Kopf-/Fußzeilen ───────────────────────────────────
\usepackage{scrlayer-scrpage}   % KOMA-kompatible Kopf-/Fußzeilen

% ── Farben ───────────────────────────────────────────────────────────
\usepackage[dvipsnames, svgnames, x11names]{xcolor}

% ── Grafiken & Abbildungen ────────────────────────────────────────────
\usepackage{graphicx}
\graphicspath{{figures/}}
\usepackage{caption}
\usepackage{subcaption}         % Teil-Abbildungen (a), (b) …
\usepackage{float}              % [H]-Platzierung
\usepackage{wrapfig}            % Text umfließt Abbildung

% ── Tabellen ─────────────────────────────────────────────────────────
\usepackage{booktabs}           % \toprule, \midrule, \bottomrule
\usepackage{tabularx}           % Spalten mit variabler Breite (X)
\usepackage{longtable}          % Mehrseitige Tabellen
\usepackage{array}              % Erweiterte Spaltentypen
\usepackage{multirow}           % Zellen über mehrere Zeilen

% ── Quellcode-Listings ───────────────────────────────────────────────
% Voraussetzung: Python + Pygments ("pip install Pygments")
% Kompilierung mit lualatex -shell-escape
\usepackage{minted}

% ── Mathematik & Einheiten ───────────────────────────────────────────
\usepackage{amsmath}
\usepackage{amssymb}
\usepackage[
  locale              = DE,
  separate-uncertainty,
  per-mode            = symbol,
]{siunitx}

% ── Abkürzungsverzeichnis ─────────────────────────────────────────────
\usepackage{acro}

% ── Aufzählungen ─────────────────────────────────────────────────────
\usepackage{enumitem}

% ── Farbige Hinweis-Boxen ─────────────────────────────────────────────
\usepackage[most]{tcolorbox}

% ── Diagramme & Plots ─────────────────────────────────────────────────
\usepackage{tikz}
\usetikzlibrary{
  arrows.meta,
  positioning,
  shapes.geometric,
  fit,
  calc,
  decorations.pathreplacing,
  backgrounds,
}
\usepackage{pgfplots}
\pgfplotsset{compat=1.18}

% ── Literaturverzeichnis ──────────────────────────────────────────────
\usepackage[
  style      = ieee,      % IEEE-Zitierstil (alternativ: authoryear, apa)
  backend    = biber,     % Moderner BibTeX-Ersatz mit Unicode-Unterstützung
  sortlocale = de_DE,
  urldate    = long,
  dashed     = false,
  maxnames   = 6,
  minnames   = 1,
]{biblatex}
\addbibresource{references.bib}

% ── Links & PDF-Metadaten (IMMER als letztes Paket!) ─────────────────
\usepackage{hyperref}
\usepackage[ngerman]{cleveref}  % Intelligente Querverweise (\cref, \Cref)
